\documentclass[12pt]{article}
\usepackage{amsmath,amssymb}

\begin{document}
\section*{{2019 AIME II Problems/Problem 8}}
Source: AoPS Wiki

\subsection*{Solution}
We have $\frac{1+\sqrt{3}i}{2} = \omega$ where $\omega = e^{\frac{i\pi}{3}}$ is a primitive 6th root of unity. Then we have \begin{align*} f(\omega) &= a\omega^{2018} + b\omega^{2017} + c\omega^{2016}\\ &= a\omega^2 + b\omega + c \end{align*} We wish to find $f(1) = a+b+c$ . We first look at the real parts. As $\text{Re}(\omega^2) = -\frac{1}{2}$ and $\text{Re}(\omega) = \frac{1}{2}$ , we have $-\frac{1}{2}a + \frac{1}{2}b + c = 2015 \implies -a+b + 2c = 4030$ . Looking at imaginary parts, we have $\text{Im}(\omega^2) = \text{Im}(\omega) = \frac{\sqrt{3}}{2}$ , so $\frac{\sqrt{3}}{2}(a+b) = 2019\sqrt{3} \implies a+b = 4038$ . As $a$ and $b$ do not exceed 2019, we must have $a = 2019$ and $b = 2019$ . Then $c = \frac{4030}{2} = 2015$ , so $f(1) = 4038 + 2015 = 6053 \implies f(1) \pmod{1000} = \boxed{053}$ . -scrabbler94

\end{document}
